\section{Berthing manipulator: kinematics and kineto-statics}
\label{sec:robotics}

\subsection{From docking to berthing}

Section~\ref{sec:proximity} delivered the chaser to the docking port with a
residual velocity of zero. In practice the last metres of a crewed vehicle's
arrival are frequently handled by a manipulator rather than by thrusters:
berthing removes the closing-rate risk, avoids plume impingement on the target,
and allows the capture to be aborted at any instant. This part characterises a
five-revolute arm sized for that job.

\subsection{Mobility}

The arm has five revolute joints connecting six bodies including the base.
Gr\"ubler's criterion in three dimensions gives

\begin{equation}
\text{DOF} = 6(N - 1 - J) + \sum_{i=1}^{J} f_i = 6(6-1-5) + 5 = 5,
\label{eq:grubler}
\end{equation}

so the arm alone has five degrees of freedom. Mounted on a free-flying base the
system has $5 + 6 = 11$. Three operating modes follow, and they are genuinely
different control problems: \emph{free-flying}, where the base pose is actively
held by thrusters; \emph{attitude-controlled}, where only orientation is held and
the base translates; and \emph{free-floating}, where nothing is actuated at the
base and total linear and angular momentum are conserved, so every arm motion
disturbs the carrier. This report assumes the free-flying case, which makes the
base pose a boundary condition rather than a state.

\subsection{Geometry and inertia}

Links are modelled as homogeneous cylinders of radius $R_c = 0.15\unit{m}$ with
the local $z$ axis along the link, so the centre of mass is at mid-length and

\begin{equation}
I_{zz} = \tfrac12 m R_c^2, \qquad
I_{xx} = I_{yy} = \tfrac{1}{12} m\left(3R_c^2 + L^2\right).
\label{eq:inertia}
\end{equation}

\begin{table}[htbp]
\centering
\caption{Link parameters. Total mass $\mArmMass\unit{kg}$, fully extended reach
$\mReach\unit{m}$. Joint axes are given in the local joint frame before the
joint rotation is applied.}
\label{tab:links}
\small
\begin{tabular}{crrrrc}
\toprule
Link & $m$ [kg] & $L$ [m] & $I_{xx}=I_{yy}$ [kg\,m$^2$] & $I_{zz}$ [kg\,m$^2$] & axis \\
\midrule
1 & $30$  & $0.60$ & $1.07$  & $0.34$ & $z$ \\
2 & $100$ & $3.20$ & $85.90$ & $1.13$ & $y$ \\
3 & $85$  & $2.80$ & $56.01$ & $0.96$ & $y$ \\
4 & $20$  & $0.60$ & $0.71$  & $0.23$ & $y$ \\
5 & $15$  & $0.50$ & $0.40$  & $0.17$ & $x$ \\
\bottomrule
\end{tabular}
\end{table}

\subsection{Forward kinematics}

Frames are placed at the link centres of mass with $z$ along the link, and
offsets are pure $z$ translations: $g_i = L_i/2$ from joint $i$ to the centre of
mass of link $i$, and $b_i = L_i/2$ from that centre of mass to joint $i+1$. The
chain is

\begin{equation}
T_i = T_{i-1}\, \mathrm{Trans}(b_{i-1})\, \mathrm{Rot}(\mathbf e_i, \theta_i)\, \mathrm{Trans}(g_i),
\label{eq:chain}
\end{equation}

with the rotation built from the Euler--Rodrigues formula

\begin{equation}
R(\mathbf e, \theta) = \cos\theta\, I + (1-\cos\theta)\,\mathbf e\mathbf e^\top + \sin\theta\,[\mathbf e\times].
\label{eq:rodrigues}
\end{equation}

The end-effector frame is one further half-link beyond the last centre of mass,
at the physical tip.

For the acceptance pose $\boldsymbol\theta = [0, 45, 0, 60, 0]^\circ$ the chain
returns a rotation of exactly $R_y(105^\circ)$ -- joints 2, 3 and 4 share the
$y$ axis, so only their sum matters -- and a tip position of
$[5.305, 0, 4.558]^\top\unit{m}$. Both reproduce the reference pose to
$6\times10^{-5}$.

\begin{figure}[htbp]
\centering
\includegraphics[width=\textwidth]{fig14_fkine_verification.png}
\caption{Forward kinematics at the acceptance pose. Left: three-dimensional
rendering with the end-effector triad. Right: $xz$ projection with the joint
origins labelled. Reach at this pose is $6.99\unit{m}$ of the $\mReach\unit{m}$
maximum.}
\label{fig:fk}
\end{figure}

\subsection{Twist propagation and kineto-static duality}

Velocities are propagated with a system Jacobian mapping the eleven generalised
velocities to the six body twists,

\begin{equation}
\mathbf t = N\,\mathbf u, \qquad
\mathbf u = \begin{bmatrix}\boldsymbol\omega_0 \\ \dot{\mathbf r}_0 \\ \dot{\mathbf q}\end{bmatrix} \in \mathbb{R}^{11},
\qquad \mathbf t \in \mathbb{R}^{36},
\qquad N = N_l N_d .
\label{eq:N}
\end{equation}

Here $N_d = \mathrm{blkdiag}(I_6, \mathbf p_1, \dots, \mathbf p_5)$ injects each
joint rate into its own body with
$\mathbf p_i = [\mathbf e_i;\ \mathbf e_i\times\mathbf g_i]$ in inertial axes,
and $N_l$ is block lower triangular with unit diagonal and rigid-body transports

\begin{equation}
B_{ij} = \begin{bmatrix} I_3 & 0 \\ [(\mathbf r_j - \mathbf r_i)\times] & I_3 \end{bmatrix}.
\label{eq:transport}
\end{equation}

The sign is the part worth checking: $\mathbf v_i = \mathbf v_j + \boldsymbol\omega_j\times(\mathbf r_i - \mathbf r_j)$
and $\boldsymbol\omega_j\times\mathbf d = -[\mathbf d\times]\boldsymbol\omega_j$,
hence $[(\mathbf r_j - \mathbf r_i)\times]$ and not the reverse. Getting it
backwards produces a base reaction whose force is right and whose moment is
wrong -- a failure mode that no plot reveals and that the tests of
Section~\ref{sec:robotics:verif} catch immediately.

In micro-gravity with a single external wrench at the end effector, static
equilibrium gives the generalised forces by duality,

\begin{equation}
\boldsymbol\tau = N^\top \mathbf w,
\label{eq:duality}
\end{equation}

where $\mathbf w$ stacks the body wrenches ordered $[\text{moment};\text{force}]$
to match the $[\boldsymbol\omega; \mathbf v]$ twist ordering. Only the last body
is loaded, with the contact wrench transported from the tip to the centre of
mass of link 5. The actuators and the attitude control system must supply
$-\boldsymbol\tau$.

\subsection{Load cases}

Two contact wrenches are applied at three postures: a purely axial
$\mathbf F = [0,0,-100]^\top\unit{N}$ representing a nominal berthing push, and a
misaligned $\mathbf F = [15,-10,-80]^\top\unit{N}$ with
$\mathbf M = [12,8,-5]^\top\unit{N\,m}$ representing an off-nominal capture.

\begin{table}[htbp]
\centering
\caption{Joint torques [N\,m] under the $100\unit{N}$ axial berthing load. The
base reaction force is $100.00\unit{N}$ in every case, as a static chain loaded
by a single external force requires.}
\label{tab:torque1}
\small
\begin{tabular}{lrrrrrr}
\toprule
Posture & $\tau_1$ & $\tau_2$ & $\tau_3$ & $\tau_4$ & $\tau_5$ & $\lVert\mathbf n_{\text{base}}\rVert$ \\
\midrule
Nearly extended & $0.00$   & $79.97$   & $24.40$   & $0.00$  & $0.00$ & $79.97$ \\
Folded          & $0.00$   & $396.47$  & $95.77$   & $0.00$  & $0.00$ & $435.03$ \\
General spatial & $131.18$ & $-339.41$ & $-194.40$ & $0.00$  & $0.00$ & $486.30$ \\
\bottomrule
\end{tabular}
\end{table}

\begin{table}[htbp]
\centering
\caption{Joint torques [N\,m] under the misaligned wrench. Base reaction force
$82.01\unit{N}$ throughout. Note that $\tau_1$ and $\tau_5$, identically zero in
the axial case, are now loaded.}
\label{tab:torque2}
\small
\begin{tabular}{lrrrrrr}
\toprule
Posture & $\tau_1$ & $\tau_2$ & $\tau_3$ & $\tau_4$ & $\tau_5$ & $\lVert\mathbf n_{\text{base}}\rVert$ \\
\midrule
Nearly extended & $-13.00$ & $177.59$  & $85.86$   & $24.50$ & $17.00$ & $206.88$ \\
Folded          & $-9.54$  & $364.72$  & $140.58$  & $24.50$ & $17.00$ & $407.65$ \\
General spatial & $62.26$  & $-193.62$ & $-112.42$ & $25.16$ & $17.00$ & $279.64$ \\
\bottomrule
\end{tabular}
\end{table}

Four observations follow, and they are the engineering content of the tables.

\begin{itemize}
\item With the arm nearly extended and the load axial, the force runs along the
      links. The $y$-axis joints see a short moment arm and the shoulder carries
      only $80\unit{N\,m}$.
\item Folding the arm turns the same $100\unit{N}$ into a long transverse lever
      and the shoulder torque rises to $\mTauMax\unit{N\,m}$, the largest value
      in the study. Actuator sizing is driven by posture, not by load magnitude.
\item The axial load leaves the base turret ($\tau_1$) and the wrist roll
      ($\tau_5$) at exactly zero, because their axes are parallel to the force or
      pass through its line of action. The misaligned wrench activates both.
      A capture mechanism designed only for the nominal case would under-size
      precisely these two actuators.
\item Masses appear nowhere in either table. In micro-gravity the static problem
      is mass-free; the inertias of Table~\ref{tab:links} would enter the mass
      matrix and Coriolis terms of the free-floating dynamics, which this report
      does not integrate.
\end{itemize}

\begin{figure}[htbp]
\centering
\includegraphics[width=\textwidth]{fig15_three_postures.png}
\caption{The three study postures under the $100\unit{N}$ axial berthing load
(arrow). Shoulder torque rises fivefold from the extended to the folded
configuration for an identical applied force.}
\label{fig:postures}
\end{figure}

\subsection{Verification}
\label{sec:robotics:verif}

Two statements must hold for a static chain loaded by a single external wrench,
and both are tested at every posture and load case:

\begin{enumerate}[label=(\roman*)]
\item the force part of the base reaction equals the applied force exactly, and
      the moment part equals $\mathbf M_{\text{tip}} + (\mathbf r_{\text{tip}} - \mathbf r_{\text{base}})\times\mathbf F$;
\item the joint torques obtained from the $36\times11$ twist-propagation route
      equal those obtained from the $6\times5$ end-effector Jacobian,
      $\boldsymbol\tau_m = J^\top[\mathbf f; \mathbf n]$.
\end{enumerate}

The worst residuals are $2\times10^{-14}\unit{N}$ and $\mDualRes\unit{N\,m}$.
Two entirely different pieces of algebra agreeing to machine precision is
stronger evidence than either one looking plausible on its own.
